\section{Uniform continuity}
\label{sec:uniform-continuity}
Uniform continuity on a set $S$ means that, for each $\varepsilon>0$,
one positive $\delta$ works for \emph{every} pair $x,y\in S$:
\[
 \abs{x-y}<\delta\quad\Longrightarrow\quad
 \abs{f(x)-f(y)}<\varepsilon.
\]
The choice may depend on $\varepsilon$ and on the fixed set $S$,
but not on the chosen points. Restricting a uniformly continuous
function to a smaller set preserves this implication, since it
reduces the number of pairs to check.

We use Theorem~19.2 for continuous functions on closed bounded
intervals and Theorem~19.6 for continuous interval functions with
bounded derivative. A bound $\abs{f'}\le M$ with $M>0$ gives
$\abs{f(x)-f(y)}\le M\abs{x-y}$ by the Mean Value theorem, so
$\delta=\varepsilon/M$ works. Derivatives are used only where
the stated interval and differentiability hypotheses are satisfied.

\exercise{19.1}{Seven tests for uniform continuity}
\goal Decide uniform continuity on each specified domain, and justify
each answer. In part (e) the function is $1/x^3$, and in part (f)
the sine argument is $1/x^2$.

\step{(a) The continuous expression on $[0,\pi]$.}
For
\[
 f(x)=x^{17}\sin x-e^x\cos(3x),
\]
the power, sine, exponential, and cosine functions are continuous.
Composition gives continuity of $\cos(3x)$, multiplication gives
continuity of both products, and subtraction gives continuity of
$f$ on $\R$. In particular it is continuous on the closed bounded
interval $[0,\pi]$. Theorem~19.2 makes $f$ uniformly continuous
there. No estimate involving its complicated derivative is needed.

\step{(b) The cube on $[0,1]$.}
The polynomial $x^3$ is continuous on $\R$, hence on $[0,1]$.
Theorem~19.2 therefore proves uniform continuity on $[0,1]$.
One can also see a point-independent estimate directly: if
$0\le x,y\le1$, then
\[
 \abs{x^3-y^3}
 =\abs{x-y}\abs{x^2+xy+y^2}\le3\abs{x-y}.
\]
Thus $\delta=\varepsilon/3$ works.

\step{(c) The cube on $(0,1)$.}
The function in (c) is the restriction of the uniformly continuous
function in (b). The same $\delta$ therefore works for all pairs
in $(0,1)$. A domain need not be closed for its function to be
uniformly continuous.

\step{(d) The cube on all of $\R$.}
For $n\in\N$, choose
\[
 x_n=n,\qquad y_n=n+\frac1{n^2}.
\]
Their distance $\abs{x_n-y_n}=1/n^2$ tends to zero, but expansion
of the cube gives
\[
 y_n^3-x_n^3
 =3n^2\frac1{n^2}+3n\frac1{n^4}+\frac1{n^6}
 =3+\frac3{n^3}+\frac1{n^6}\ge3.
\]
Fix $\varepsilon_0=1$. For any proposed $\delta>0$, choose
$n$ so large that $1/n^2<\delta$. The corresponding pair is
closer than $\delta$, but its function values differ by at least
$3>\varepsilon_0$. This violates Definition~19.1. The cube is
not uniformly continuous on $\R$.

\step{(e) The reciprocal cube on $(0,1]$.}
The sequence $x_n=1/n$ lies in $(0,1]$ and converges to zero,
so it is a Cauchy sequence. For $f(x)=1/x^3$ its image is
$f(x_n)=n^3$, which is unbounded. Every real Cauchy sequence is
bounded, so this image is not a Cauchy sequence. If $f$ were
uniformly continuous, Theorem~19.4 would make the image of every
Cauchy sequence a Cauchy sequence. This contradiction proves
that $1/x^3$ is not uniformly continuous on $(0,1]$.

\step{(f) The rapidly oscillating sine on $(0,1]$.}
Let $f(x)=\sin(1/x^2)$ and choose
\[
 x_n=\frac1{\sqrt{2\pi n+\pi/2}},\qquad
 y_n=\frac1{\sqrt{2\pi n+3\pi/2}}.
\]
The denominators exceed $1$, so both sequences lie in $(0,1]$.
Both tend to zero, whence
$\abs{x_n-y_n}\le x_n+y_n\to0$. On the other hand,
\[
 f(x_n)=\sin(2\pi n+\pi/2)=1,\qquad
 f(y_n)=\sin(2\pi n+3\pi/2)=-1.
\]
Their images always differ by $2$. Taking $\varepsilon_0=1$,
every proposed $\delta>0$ is defeated by a sufficiently late
pair with $\abs{x_n-y_n}<\delta$. Thus $f$ is not uniformly
continuous on this interval.

\step{(g) The damped oscillation on $(0,1]$.}
For $x>0$ let $f(x)=x^2\sin(1/x)$, and extend the formula by
$F(0)=0$, $F(x)=f(x)$ for $0<x\le1$. For positive $x$ it is a
product and composition of continuous functions. At zero,
\[
 \abs{F(x)-F(0)}=\abs{x^2\sin(1/x)}\le x^2.
\]
Given $\varepsilon>0$, the choice $\delta=\sqrt\varepsilon$
makes this error less than $\varepsilon$ whenever
$0\le x<\delta$. Thus $F$ is continuous on $[0,1]$.
Theorem~19.2 makes $F$ uniformly continuous there, and restriction
proves uniform continuity of $f$ on $(0,1]$.
\answer{Parts (a), (b), (c), and (g) are uniformly continuous on
their specified sets. Parts (d), (e), and (f) are not.}

\exercise{19.2}{Three direct epsilon--delta proofs}
\goal Verify Definition~19.1 directly for each function, choosing
$\delta$ independently of the pair of points.

\step{(a) The affine function $3x+11$ on $\R$.}
Let $\varepsilon>0$ and set $\delta=\varepsilon/3>0$.
For arbitrary $x,y\in\R$ with $\abs{x-y}<\delta$, cancellation
of the constants gives
\[
 \abs{(3x+11)-(3y+11)}
 =3\abs{x-y}<3\delta=\varepsilon.
\]
This one $\delta$ works throughout the unbounded domain.

\step{(b) The square on $[0,3]$.}
For $x,y\in[0,3]$, one has $0\le x+y\le6$. Factor the
difference of squares to get
\[
 \abs{x^2-y^2}=\abs{x-y}\abs{x+y}\le6\abs{x-y}.
\]
For any $\varepsilon>0$, take $\delta=\varepsilon/6$.
Then $\abs{x-y}<\delta$ makes the displayed error less than
$6\delta=\varepsilon$, for every such pair.

\step{(c) The reciprocal on $[1/2,+\infty)$.}
If $x,y\ge1/2$, both denominators are positive and
$xy\ge1/4$. Consequently
\[
 \abs{\frac1x-\frac1y}
 =\frac{\abs{y-x}}{xy}\le4\abs{x-y}.
\]
Choose $\delta=\varepsilon/4$. For every pair in the domain
with distance less than $\delta$, the function-value distance is
less than $4\delta=\varepsilon$.
\answer{The respective choices $\varepsilon/3$, $\varepsilon/6$,
and $\varepsilon/4$ establish all three uniform-continuity claims
directly from the definition.}

\exercise{19.3}{Direct proofs for two rational functions}
\goal Repeat the direct epsilon--delta verification for
$x/(x+1)$ on $[0,2]$ and $5x/(2x-1)$ on $[1,+\infty)$.

\step{(a) Subtract with the denominator kept explicit.}
For $x,y\in[0,2]$, both $x+1$ and $y+1$ are at least $1$.
Common denominators and cancellation give
\[
 \begin{aligned}
 \abs{\frac{x}{x+1}-\frac{y}{y+1}}
 &=\frac{\abs{x(y+1)-y(x+1)}}{(x+1)(y+1)}\\
 &=\frac{\abs{xy+x-xy-y}}{(x+1)(y+1)}
 =\frac{\abs{x-y}}{(x+1)(y+1)}\le\abs{x-y}.
 \end{aligned}
\]
Given $\varepsilon>0$, set $\delta=\varepsilon$. The final bound
then becomes less than $\varepsilon$ for every pair with
$\abs{x-y}<\delta$.

\step{(b) Do the same calculation on the unbounded interval.}
For $x,y\ge1$, one has $2x-1\ge1$ and $2y-1\ge1$.
Thus
\[
 \begin{aligned}
 \abs{\frac{5x}{2x-1}-\frac{5y}{2y-1}}
 &=\frac{\abs{5x(2y-1)-5y(2x-1)}}{(2x-1)(2y-1)}\\
 &=\frac{5\abs{-x+y}}{(2x-1)(2y-1)}
 \le5\abs{x-y}.
 \end{aligned}
\]
Set $\delta=\varepsilon/5$. The error is then less than
$5\delta=\varepsilon$ whenever $\abs{x-y}<\delta$.
\answer{Both functions are uniformly continuous. The direct
choices are $\delta=\varepsilon$ in (a) and
$\delta=\varepsilon/5$ in (b).}

\exercise{19.4}{A uniformly continuous function on a bounded set is bounded}
\goal Prove the boundedness assertion in (a), without assuming
the domain is an interval or is closed, and apply it to part (b).

\step{(a), preparation: account for the empty domain.}
If $S=\varnothing$, every bound is a bound on its function values,
so the conclusion is immediate. Suppose henceforth that $S$ is
nonempty and bounded, and that $f$ is uniformly continuous on $S$.

\step{(a), contradiction: select values that are too large.}
If $f$ were unbounded, then for every integer $n\ge1$ some
$x_n\in S$ would satisfy $\abs{f(x_n)}>n$. Otherwise a bound
on all $\abs{f(x)}$ would exist, which is exactly boundedness
of $f$. Since $S$ is bounded, the real sequence $(x_n)$ is
bounded. Bolzano--Weierstrass (Theorem~11.5) supplies a subsequence
$x_{n_j}$ converging to some real number $x_0$.
It is not necessary that $x_0$ belong to $S$.

Every convergent real sequence is Cauchy, so $(x_{n_j})$ is a
Cauchy sequence of points in $S$. Theorem~19.4 then makes
$(f(x_{n_j}))$ a Cauchy sequence, hence a bounded sequence.
However, its construction gives
\[
 \abs{f(x_{n_j})}>n_j\ge j.
\]
Here $n_j\ge j$ because the selected indices are strictly
increasing positive integers. The displayed lower bound makes
the image sequence unbounded, a contradiction. Thus $f$ is bounded.

\step{(b) Apply the assertion to the reciprocal square.}
The set $(0,1)$ is bounded. For $n\ge1$ the points
$x_n=1/(n+1)$ lie in $(0,1)$, while
\[
 \frac1{x_n^2}=(n+1)^2\longrightarrow+\infty.
\]
The function $1/x^2$ is therefore unbounded on this bounded
domain. If it were uniformly continuous, part (a) would make
it bounded. Hence it is not uniformly continuous on $(0,1)$.
\answer{Uniform continuity implies boundedness on every bounded
real domain, including domains with missing limit points.
This excludes uniform continuity of $1/x^2$ on $(0,1)$.}

\exercise{19.5}{Six domains for trigonometric and reciprocal functions}
\goal Determine which of the six specified restrictions are
uniformly continuous, with the domain included in each argument.

\step{(a) The tangent on $[0,\pi/4]$.}
Cosine has no zero in this closed interval. Thus
$\tan x=\sin x/\cos x$ is continuous there, and Theorem~19.2
gives uniform continuity.

\step{(b) The tangent on $[0,\pi/2)$.}
This domain is bounded, but tangent is unbounded on it.
For example, let $x_n=\pi/2-1/n$, which lies in the domain.
Then
\[
 \sin x_n\to1,\qquad \cos x_n\to0,
 \qquad \cos x_n>0.
\]
For any $M>0$, the first convergence gives $\sin x_n>1/2$
eventually, and the second gives $0<\cos x_n<1/(2M)$
eventually. At indices satisfying both estimates,
$\tan x_n>M$. Thus tangent is unbounded and cannot be
uniformly continuous on this bounded set by Exercise~19.4(a).

\step{(c) The function $(\sin^2x)/x$ on $(0,\pi]$.}
We need the elementary estimate $\abs{\sin t}\le\abs{t}$.
It follows from the Mean Value theorem applied to sine between
$0$ and $t$: its derivative is cosine, with absolute value at
most $1$. Exercise~19.8 supplies the more general version.
For $x>0$ this gives
\[
 0\le\frac{\sin^2x}{x}\le\frac{x^2}{x}=x.
\]
Extend the function to $[0,\pi]$ by assigning value $0$ at zero.
The displayed bound proves continuity at zero: if
$0\le x<\varepsilon$, its value has absolute value less than
$\varepsilon$. Away from zero the quotient is continuous.
The extension is therefore continuous on $[0,\pi]$, hence
uniformly continuous by Theorem~19.2. Its restriction to
$(0,\pi]$ is uniformly continuous.

\step{(d) The function $1/(x-3)$ on $(0,3)$.}
For $x_n=3-1/n$, one has $2\le x_n<3$, so $x_n$ is in
the domain, and
\[
 \frac1{x_n-3}=-n.
\]
The function is unbounded on the bounded set $(0,3)$.
Exercise~19.4(a) therefore excludes uniform continuity.

\step{(e) The same function on $(3,+\infty)$.}
This domain is unbounded, so unboundedness of the function alone
would not settle uniform continuity. Instead take
$x_n=3+1/n$. This converges to $3$ and is therefore a Cauchy
sequence of points in the domain. Its image is
\[
 \frac1{x_n-3}=n,
\]
which is unbounded and hence is not Cauchy. The Cauchy-sequence
preservation theorem (Theorem~19.4) excludes uniform continuity.

\step{(f) The same function on $(4,+\infty)$.}
For $x,y>4$, both $x-3$ and $y-3$ exceed $1$.
Thus
\[
 \abs{\frac1{x-3}-\frac1{y-3}}
 =\frac{\abs{x-y}}{(x-3)(y-3)}\le\abs{x-y}.
\]
The choice $\delta=\varepsilon$ works for every pair in this
domain. Hence the function is uniformly continuous here.
\answer{Parts (a), (c), and (f) are uniformly continuous;
parts (b), (d), and (e) are not.}

\exercise{19.6}{Uniform continuity without a bounded derivative}
\goal For $f(x)=\sqrt x$, show the derivative is unbounded on
$(0,1]$ while $f$ is uniformly continuous there, and prove uniform
continuity on $[1,+\infty)$.

\step{(a), derivative: exhibit arbitrarily large values.}
For $x>0$, elementary differentiation gives
$f'(x)=1/(2\sqrt x)$. The points $x_n=1/n^2$ belong to
$(0,1]$ and satisfy
\[
 f'(x_n)=\frac{1}{2(1/n)}=\frac n2\longrightarrow+\infty.
\]
Thus the derivative is unbounded on this interval.

\step{(a), the function: use the square-root difference estimate.}
For $x\ge y\ge0$, both $\sqrt x-\sqrt y$ and
$\sqrt x+\sqrt y$ are nonnegative, and the latter is at least
the former. Therefore
\[
 x-y=(\sqrt x-\sqrt y)(\sqrt x+\sqrt y)
 \ge(\sqrt x-\sqrt y)^2.
\]
Taking nonnegative square roots yields
$\abs{\sqrt x-\sqrt y}\le\sqrt{\abs{x-y}}$.
If $y>x$, interchange the two points to obtain the same estimate.
For $\varepsilon>0$, choose $\delta=\varepsilon^2$.
Whenever $x,y\in(0,1]$ and $\abs{x-y}<\delta$, it follows that
\[
 \abs{\sqrt x-\sqrt y}\le\sqrt{\abs{x-y}}
 <\sqrt\delta=\varepsilon.
\]
This establishes uniform continuity despite the unbounded
derivative. Theorem~19.6 gives a sufficient condition, not a
necessary condition, so there is no contradiction.

\step{(b) Rationalize on the domain bounded away from zero.}
For $x,y\ge1$, one has $\sqrt x+\sqrt y\ge2$. Hence
\[
 \abs{\sqrt x-\sqrt y}
 =\frac{\abs{x-y}}{\sqrt x+\sqrt y}\le\frac12\abs{x-y}.
\]
Taking $\delta=2\varepsilon$ proves uniform continuity on
$[1,+\infty)$ directly.
\answer{The square root is uniformly continuous on both specified
domains. Its derivative is nevertheless unbounded near zero.}

\exercise{19.7}{Joining a compact initial interval to a uniform tail}
\goal Prove part (a)'s extension from a tail to the whole
nonnegative half-line, and apply it to square root in part (b).

\step{(a), preparation: choose overlapping intervals.}
Let $f$ be continuous on $[0,+\infty)$ and uniformly continuous
on $[k,+\infty)$ as in the hypothesis. We only need the part
of that tail contained in its domain. Put $K=\max\{k,0\}$;
restriction gives uniform continuity on $[K,+\infty)$.
If $k\le0$ and the stated tail hypothesis is meaningful, it
already includes $[0,+\infty)$, but the argument below also works.

The restriction to $[0,K+1]$ is continuous on a closed bounded
interval, hence uniformly continuous by Theorem~19.2. These two
intervals overlap on $[K,K+1]$. The overlap will handle pairs
that are close but lie on opposite sides of the cutoff.

\step{(a), verification: use one delta for every pair.}
Fix $\varepsilon>0$. Choose $\delta_0>0$ for uniform continuity
on $[0,K+1]$, and $\delta_1>0$ for uniform continuity on
$[K,+\infty)$, both for this same tolerance $\varepsilon$.
Set
\[
 \delta=\min\{\delta_0,\delta_1,1\}>0.
\]
Take $x,y\ge0$ with $\abs{x-y}<\delta$. Interchanging them
if necessary, assume $x\le y$.
If $y\le K+1$, both points lie in the first interval and
their distance is less than $\delta_0$, so their function values
differ by less than $\varepsilon$.
If $y>K+1$, then $y-x<1$, which rearranges to
$x>y-1>K$. Thus both points lie in the tail, and their distance
is less than $\delta_1$, so the same conclusion holds.
These two cases exhaust every pair. Thus $f$ is uniformly
continuous on the whole half-line.

\step{(b) Apply part (a) to square root.}
The function $\sqrt x$ is continuous on $[0,+\infty)$ by
Exercise~17.4. Exercise~19.6(b) gives uniform continuity on
$[1,+\infty)$. Part (a), with $k=1$, now proves uniform
continuity on $[0,+\infty)$.
\answer{A continuous function on the half-line is uniformly
continuous whenever it is uniformly continuous on a tail.
The square root is one such function.}

\exercise{19.8}{The sine difference estimate}
\goal Use the Mean Value theorem to prove
$\abs{\sin x-\sin y}\le\abs{x-y}$, then deduce uniform
continuity of sine on $\R$.

\step{(a) Apply the theorem on the interval between the points.}
If $x=y$, both sides are zero. Suppose $x<y$. Sine is
continuous on $[x,y]$ and differentiable on $(x,y)$, so the
Mean Value theorem gives some $c\in(x,y)$ with
\[
 \sin y-\sin x=(\cos c)(y-x).
\]
Taking absolute values and using $\abs{\cos c}\le1$ gives
\[
 \abs{\sin y-\sin x}
 =\abs{\cos c}\abs{y-x}\le\abs{y-x}.
\]
If $x>y$, interchange the points; the absolute values make
the desired statement unchanged. This proves the estimate for
all pairs of real numbers.

\step{(b) Read off a uniform choice.}
Given $\varepsilon>0$, let $\delta=\varepsilon$.
For all $x,y\in\R$ with $\abs{x-y}<\delta$, part (a) gives
\[
 \abs{\sin x-\sin y}\le\abs{x-y}<\delta=\varepsilon.
\]
The choice is independent of both points.
\answer{Sine is uniformly continuous on $\R$, with the global
difference bound claimed in (a). In particular,
$\abs{\sin t}\le\abs t$ follows by setting $y=0$.}

\exercise{19.9}{A damped oscillation on the whole real line}
\goal For $f(0)=0$ and $f(x)=x\sin(1/x)$ when $x\ne0$,
verify continuity, uniform continuity on every bounded subset,
and decide uniform continuity on $\R$.

\step{(a) Check the origin as well as the ordinary formula.}
At every nonzero point, the reciprocal, sine composition,
and product are continuous. At zero,
\[
 \abs{f(x)-f(0)}\le\abs x
\]
for $x\ne0$, since $\abs{\sin(1/x)}\le1$; the inequality
also holds at zero. Thus $\delta=\varepsilon$ proves continuity
there. The function is continuous on all of $\R$.

\step{(b) Enclose any bounded set in a closed interval.}
Let $S\subseteq\R$ be bounded. Choose $R>0$ so that
$S\subseteq[-R,R]$. The function is continuous on this
closed bounded interval, hence uniformly continuous there by
Theorem~19.2. Restriction gives uniform continuity on $S$,
even if $S$ is not closed or is empty.

\step{(c), tails: bound the derivative away from the origin.}
For $x\ne0$, the product and chain rules give
\[
 f'(x)=\sin(1/x)-\frac{\cos(1/x)}{x}.
\]
If $\abs x\ge1$, then
\[
 \abs{f'(x)}\le1+\frac1{\abs x}\le2.
\]
Theorem~19.6 therefore gives uniform continuity on each of the
intervals $[1,+\infty)$ and $(-\infty,-1]$. More explicitly,
the Mean Value theorem gives the bound
$\abs{f(x)-f(y)}\le2\abs{x-y}$ whenever both points lie
in either one of these intervals.

\step{(c), combine the tails and the central interval.}
The central restriction to $[-2,2]$ is uniformly continuous
by part (b). For a fixed $\varepsilon>0$, choose its uniform
parameter $\delta_0>0$ and put
\[
 \delta=\min\{\delta_0,\varepsilon/2,1\}>0.
\]
Take $x\le y$ with $y-x<\delta$.
If both points lie in $[-2,2]$, the central estimate gives
$\abs{f(x)-f(y)}<\varepsilon$.
If $y>2$, then $x>y-1>1$, so both lie in the positive tail;
the derivative bound gives
$\abs{f(x)-f(y)}\le2\abs{x-y}<2\delta\le\varepsilon$.
If $x<-2$, then $y<x+1<-1$, so both lie in the negative
tail, and the same bound applies.
These cases cover every pair: a pair outside the central square
has either $y>2$ or $x<-2$. Thus the function is uniformly
continuous on $\R$.

This last argument also proves a useful joining fact: a continuous
real function that is uniformly continuous on both tails
$[1,+\infty)$ and $(-\infty,-1]$ is uniformly continuous on
$\R$. In the general case, replace $\varepsilon/2$ above by
the minimum of the two tail parameters. The same interval
classification is unchanged.
\answer{The function is continuous on $\R$, uniformly continuous
on every bounded subset, and also uniformly continuous on all
of $\R$.}

\exercise{19.10}{The quadratic damping has uniform control as well}
\goal Repeat all three parts of Exercise~19.9 for
$g(0)=0$ and $g(x)=x^2\sin(1/x)$ when $x\ne0$.

\step{(a) Verify continuity everywhere.}
Away from zero, compositions and products give continuity.
At zero,
\[
 \abs{g(x)-g(0)}\le x^2.
\]
For $\varepsilon>0$, choose $\delta=\sqrt\varepsilon$.
Then $\abs{x}<\delta$ makes the displayed error less than
$\varepsilon$, so $g$ is continuous at zero as well.

\step{(b) Restrict a continuous function on a larger compact interval.}
For any bounded $S\subseteq\R$, choose $R>0$ containing it
in $[-R,R]$. Continuity on that interval implies uniform
continuity by Theorem~19.2. Restriction proves uniform continuity
on $S$.

\step{(c) Obtain a derivative bound on all of $\R$.}
For $x\ne0$, differentiation gives
\[
 g'(x)=2x\sin(1/x)-\cos(1/x).
\]
By Exercise~19.8(a), $\abs{\sin(1/x)}\le1/\abs x$.
Thus $\abs{x\sin(1/x)}\le1$ for every $x\ne0$, and
\[
 \abs{g'(x)}\le2\abs{x\sin(1/x)}+\abs{\cos(1/x)}
 \le2+1=3.
\]
The derivative at zero exists too. Indeed, for $h\ne0$,
\[
 \abs{\frac{g(h)-g(0)}{h}-0}
 =\abs{h\sin(1/h)}\le\abs h\longrightarrow0.
\]
This is precisely the difference quotient defining $g'(0)$,
so $g'(0)=0$. Consequently $g$ is differentiable on all of
$\R$ and satisfies $\abs{g'}\le3$ everywhere.

The Mean Value theorem now gives
$\abs{g(x)-g(y)}\le3\abs{x-y}$ for every pair of real
points. The choice $\delta=\varepsilon/3$ proves uniform
continuity on $\R$, also as an application of Theorem~19.6.
\answer{All three conclusions are affirmative: $g$ is continuous
on $\R$, uniformly continuous on every bounded subset, and
uniformly continuous on $\R$. Its global derivative bound is $3$.}

\exercise{19.11}{Uniform continuity of the extended sinc function}
\goal Accept the continuity on $\R$ of the extension
\[
 \widetilde h(x)=
 \begin{cases}
  \sin x/x,&x\ne0,\\
  1,&x=0,
 \end{cases}
\]
from Example~9, and prove uniform continuity on $\R$.

\step{1. Bound the derivative on both tails.}
For $x\ne0$, the quotient rule gives
\[
 \widetilde h'(x)=\frac{x\cos x-\sin x}{x^2}.
\]
If $\abs x\ge1$, the real triangle inequality yields
\[
 \abs{\widetilde h'(x)}
 \le\frac{\abs x\abs{\cos x}+\abs{\sin x}}{x^2}
 \le\frac1{\abs x}+\frac1{x^2}\le2.
\]
The function is continuous on each tail and differentiable in
its interior, so Theorem~19.6 gives uniform continuity on
$[1,+\infty)$ and on $(-\infty,-1]$.

\step{2. Join the two tails through a compact middle interval.}
The continuity assumed in the exercise gives uniform continuity
on $[-2,2]$ by Theorem~19.2. Fix $\varepsilon>0$, let
$\delta_0>0$ be its central parameter, and set
$\delta=\min\{\delta_0,\varepsilon/2,1\}$.
For a pair $x\le y$ with $y-x<\delta$, if both are in
$[-2,2]$, the central estimate works. If $y>2$, then
$x>y-1>1$ and both are in the positive tail. If $x<-2$,
then $y<x+1<-1$ and both are in the negative tail.
On either tail the Mean Value theorem and the bound $2$
give an error at most $2\abs{x-y}<\varepsilon$.
Every pair belongs to one of these cases, which proves the
uniform-continuity definition on all of $\R$.
\answer{The continuous extension of $\sin x/x$ at zero is uniformly
continuous on the whole real line. The argument uses the stipulated
continuity at zero and derivative bounds only away from zero.}

\exercise{19.12}{The running maximum of a continuous function}
\goal For a continuous $f:[a,b]\to\R$, prove that
\[
 f^*(x)=\sup\{f(t):a\le t\le x\},\qquad x\in[a,b],
\]
is increasing and continuous. Increasing here means nondecreasing,
as elsewhere in the text.

\step{1. Verify that every supremum is finite and attained.}
For each $x\in[a,b]$, the interval $[a,x]$ is nonempty,
closed, and bounded. The restriction of $f$ is continuous on it.
The extreme value theorem (Theorem~18.1) therefore supplies
some $z\in[a,x]$ with $f(z)=\max_{a\le t\le x}f(t)$.
In particular the supremum defining $f^*(x)$ is a real number
and is this maximum. When $x=a$, the interval is the singleton
$\{a\}$ and the assertion reduces to $f^*(a)=f(a)$.

\step{2. Check the order directly from set inclusion.}
If $a\le x\le y\le b$, then $[a,x]\subseteq[a,y]$.
Every function value used in the supremum for $x$ is therefore
also used for $y$. Hence $f^*(x)\le f^*(y)$. This proves
the required increasing property; it need not be strict,
as a constant $f$ demonstrates.

\step{3. Obtain one parameter from uniform continuity of $f$.}
Theorem~19.2 makes $f$ uniformly continuous on $[a,b]$.
For $\varepsilon>0$, choose $\delta>0$ such that
\[
 u,v\in[a,b],\quad\abs{u-v}<\delta
 \quad\Longrightarrow\quad\abs{f(u)-f(v)}<\varepsilon.
\]
We prove that this same $\delta$ works for $f^*$.

\step{4. Locate a point attaining the larger running maximum.}
Take $x,y\in[a,b]$ with $\abs{x-y}<\delta$.
If $x=y$, the two running maxima agree. Otherwise interchange
the points if needed so that $x<y$. By Step~1 choose
$z\in[a,y]$ with $f(z)=f^*(y)$.

If $z\le x$, then $f(z)\le f^*(x)$, since $z$ is among
the points considered in that smaller supremum. Combining this
with Step~2 gives $f^*(y)\le f^*(x)\le f^*(y)$, so the
two values agree and their difference is zero.

If $z>x$, then $0<z-x\le y-x<\delta$. The chosen parameter
therefore gives $\abs{f(z)-f(x)}<\varepsilon$.
Since $f(x)\le f^*(x)$ and $f(z)=f^*(y)$, we get
\[
 0\le f^*(y)-f^*(x)
 =f(z)-f^*(x)
 \le f(z)-f(x)
 \le\abs{f(z)-f(x)}<\varepsilon.
\]
The initial inequality again uses monotonicity from Step~2.
These two cases prove
$\abs{f^*(y)-f^*(x)}<\varepsilon$ for every eligible pair.

Thus $f^*$ is actually uniformly continuous on $[a,b]$.
In particular, fixing either point as an arbitrary $x_0$ verifies
continuity at $x_0$, including the endpoints. If $a=b$, there
is only one point and the same assertion is immediate.
\answer{The running supremum is a finite attained maximum,
is nondecreasing, and is continuous on $[a,b]$; indeed, it is
uniformly continuous there.}
